% GENERATED FILE. Edit canonical lessons, then run npm run generate:books.
% canonical-source-hash: fnv1a64:c393c5fad2b3b9f8
% canonical-lessons: BN-C81-kachhe, BN-C81-bhitore, BN-C81-baire, BN-C81-niche, BN-C81-samne

\chapter{Near, Inside, Outside, Below, In front}
\label{ch:bn-near-inside-outside-below-in-front}
\hlchaptermodality{81}

\begin{chapteropening}
\textbf{By the end of this chapter:} \emph{I can say near, inside, outside, below and in front.}

Complete the last lesson of chapter 81: I can say near, inside, outside, below and in front.
\end{chapteropening}

\section[\texorpdfstring{\bn{কাছে} (kāchhe)}{কাছে (kāchhe)}]{\bn{কাছে} (kāchhe) — near}

\label{lesson:BN-C81-kachhe}



\begin{quote}
Before the new one: say the Bengali for time, then the Bengali for midday.
\end{quote}

\subsection*{You'll want to know: \bn{কাছে}}
\textbf{\bn{কাছে}} — \emph{kāchhe} — \textquotedblleft{}near\textquotedblright{}.

\begin{cousinweb}
The origin of \textbf{\bn{কাছে}} (\emph{kāchhe}) is not settled. CDIAL has no entry for it; under Sanskrit \emph{kakṣa}, \textquotedblleft{}armpit, side\textquotedblright{}, it lists Bengali \emph{kāṁkh}, \textquotedblleft{}flank\textquotedblright{}, and says some postpositions probably grew from the sense \textquotedblleft{}side\textquotedblright{} (CDIAL 2588), but it does not trace \textbf{\bn{কাছে}} (\emph{kāchhe}) there.
\end{cousinweb}

\begin{grammarlens}[title={What you've built}]
One more word for this chapter.
\end{grammarlens}

\subsection*{Your turn}
\begin{itemize}
\raggedright
  \item \emph{Say it:} \emph{kāchhe}
  \item \emph{Say it:} \emph{kāchhe}, once more
  \item \emph{Say it:} say \emph{kāchhe}
  \item \emph{Recall:} say the Bengali for time, then the Bengali for midday, then say \emph{kāchhe} again
\end{itemize}

\begin{tcolorbox}[breakable,colback=teal!4,colframe=teal!35!black,title={Before you move on}]
What does \bn{কাছে} mean? (\textquotedblleft{}Near\textquotedblright{}.) Say it once more.
\end{tcolorbox}
\section[\texorpdfstring{\bn{ভিতরে} (bhitôre)}{ভিতরে (bhitôre)}]{\bn{ভিতরে} (bhitôre) — inside}

\label{lesson:BN-C81-bhitore}



\begin{quote}
Before the new one: say the Bengali for midday, then the Bengali for near.
\end{quote}

\subsection*{You'll want to know: \bn{ভিতরে}}
\textbf{\bn{ভিতরে}} — \emph{bhitôre} — \textquotedblleft{}inside\textquotedblright{}.

\begin{cousinweb}
\textbf{\bn{ভিতরে}} (\emph{bhitôre}) is a tadbhava, a word inherited and reshaped by sound change. It comes from \emph{*bhiyantara}, \textquotedblleft{}inner\textquotedblright{}, a form CDIAL reconstructs from its descendants (CDIAL 9504). It stands beside Sanskrit \emph{abhyantara}, \textquotedblleft{}inner\textquotedblright{}; CDIAL lists Bengali \emph{bhitar}, \textquotedblleft{}inside\textquotedblright{}.
\end{cousinweb}

\begin{grammarlens}[title={What you've built}]
One more word for this chapter.
\end{grammarlens}

\subsection*{Your turn}
\begin{itemize}
\raggedright
  \item \emph{Say it:} \emph{bhitôre}
  \item \emph{Say it:} \emph{bhitôre}, once more
  \item \emph{Say it:} say \emph{bhitôre}
  \item \emph{Recall:} say the Bengali for midday, then the Bengali for near, then say \emph{bhitôre} again
\end{itemize}

\begin{tcolorbox}[breakable,colback=teal!4,colframe=teal!35!black,title={Before you move on}]
What does \bn{ভিতরে} mean? (\textquotedblleft{}Inside\textquotedblright{}.) Say it once more.
\end{tcolorbox}
\section[\texorpdfstring{\bn{বাইরে} (bāire)}{বাইরে (bāire)}]{\bn{বাইরে} (bāire) — outside}

\label{lesson:BN-C81-baire}



\begin{quote}
Before the new one: say the Bengali for near, then the Bengali for inside.
\end{quote}

\subsection*{You'll want to know: \bn{বাইরে}}
\textbf{\bn{বাইরে}} — \emph{bāire} — \textquotedblleft{}outside\textquotedblright{}.

\begin{cousinweb}
\textbf{\bn{বাইরে}} (\emph{bāire}) is a tadbhava, a word inherited and reshaped by sound change. It comes from \emph{*bāhira}, \textquotedblleft{}external\textquotedblright{}, a form CDIAL reconstructs from its descendants (CDIAL 9226). CDIAL lists Bengali \emph{bāhir} and \emph{bār}; \emph{bāire} has the ending \emph{-e}.
\end{cousinweb}

\begin{grammarlens}[title={What you've built}]
One more word for this chapter.
\end{grammarlens}

\subsection*{Your turn}
\begin{itemize}
\raggedright
  \item \emph{Say it:} \emph{bāire}
  \item \emph{Say it:} \emph{bāire}, once more
  \item \emph{Say it:} say \emph{bāire}
  \item \emph{Recall:} say the Bengali for near, then the Bengali for inside, then say \emph{bāire} again
\end{itemize}

\begin{tcolorbox}[breakable,colback=teal!4,colframe=teal!35!black,title={Before you move on}]
What does \bn{বাইরে} mean? (\textquotedblleft{}Outside\textquotedblright{}.) Say it once more.
\end{tcolorbox}
\section[\texorpdfstring{\bn{নিচে} (niche)}{নিচে (niche)}]{\bn{নিচে} (niche) — below}

\label{lesson:BN-C81-niche}



\begin{quote}
Before the new one: say the Bengali for inside, then the Bengali for outside.
\end{quote}

\subsection*{You'll want to know: \bn{নিচে}}
\textbf{\bn{নিচে}} — \emph{niche} — \textquotedblleft{}below\textquotedblright{}.

\begin{cousinweb}
\textbf{\bn{নিচে}} (\emph{niche}) is a tadbhava, a word inherited from Sanskrit and reshaped by sound change. It comes from \emph{nīca}, \textquotedblleft{}low\textquotedblright{} (CDIAL 7540). \emph{Niche} adds the ending \emph{-e}, \textquotedblleft{}below\textquotedblright{}.
\end{cousinweb}

\begin{grammarlens}[title={What you've built}]
One more word for this chapter.
\end{grammarlens}

\subsection*{Your turn}
\begin{itemize}
\raggedright
  \item \emph{Say it:} \emph{niche}
  \item \emph{Say it:} \emph{niche}, once more
  \item \emph{Say it:} say \emph{niche}
  \item \emph{Recall:} say the Bengali for inside, then the Bengali for outside, then say \emph{niche} again
\end{itemize}

\begin{tcolorbox}[breakable,colback=teal!4,colframe=teal!35!black,title={Before you move on}]
What does \bn{নিচে} mean? (\textquotedblleft{}Below\textquotedblright{}.) Say it once more.
\end{tcolorbox}
\section[\texorpdfstring{\bn{সামনে} (sāmne)}{সামনে (sāmne)}]{\bn{সামনে} (sāmne) — in front}

\label{lesson:BN-C81-samne}



\begin{quote}
Before the new one: say the Bengali for outside, then the Bengali for below.
\end{quote}

\subsection*{You'll want to know: \bn{সামনে}}
\textbf{\bn{সামনে}} — \emph{sāmne} — \textquotedblleft{}in front\textquotedblright{}.

\begin{cousinweb}
\textbf{\bn{সামনে}} (\emph{sāmne}) goes back to Sanskrit \emph{saṁmukha}, \textquotedblleft{}facing, present\textquotedblright{}. CDIAL gives no Bengali form under that entry, but it lists Middle Bengali \emph{chāmu}, \textquotedblleft{}in front\textquotedblright{}, and Oriya \emph{sāmanā}, \textquotedblleft{}being in front\textquotedblright{}, the forms nearest to Bengali (CDIAL 12982).
\end{cousinweb}

\begin{grammarlens}[title={What you've built}]
One more word for this chapter.
\end{grammarlens}

\subsection*{Your turn}
\begin{itemize}
\raggedright
  \item \emph{Say it:} \emph{sāmne}
  \item \emph{Say it:} \emph{sāmne}, once more
  \item \emph{Say it:} say \emph{sāmne}
  \item \emph{Recall:} say the Bengali for outside, then the Bengali for below, then say \emph{sāmne} again
\end{itemize}

\begin{tcolorbox}[breakable,colback=teal!4,colframe=teal!35!black,title={Before you move on}]
What does \bn{সামনে} mean? (\textquotedblleft{}In front\textquotedblright{}.) Say it once more.
\end{tcolorbox}
